\maketitle\setcounter{page}{2}
\chapter*{摘要}
本文采用明确的反应规则与动力学假设，从产生--降解、启动子结合和质量作用定律出发，研究少组分基因调控回路所允许的行为。全文围绕一个问题展开：调控如何把简单的浓度恢复变成记忆、振荡或适应？先从反应建立表达模型，再依次加入自反馈、相互反馈和第三个调控环节，沿着每个模型的稳态与扰动响应解释这些行为的来源。双节点部分完整整理 81 种符号拓扑，进一步结合 Science（2008）的正负反馈振荡研究与 Cell（2009）的适应网络枚举，比较可调时钟、时间筛选、缓冲反馈与比例前馈。随后把连续浓度改为离散分子数，最后用乳糖操纵子讨论最小回路如何连接真实网络。重点在于区分\emph{结构所允许的机制}、\emph{特定参数下实现的动力学}与\emph{具体生物学功能}，而非追求复杂网络的完整重构。

\clearpage{\footnotesize\linespread{.93}\selectfont\setlength{\parskip}{0pt}\tableofcontents}\clearpage
